% =============================================================
% §8 Conclusion
% =============================================================
\section{Conclusion}
\label{sec:conclusion}

This paper investigated AIGC-aware reinforcement-learning
scheduling for cloud-edge LLM inference. We built a
comprehensive AIGC simulator modeling five physical
characteristics---model weight residency, KV-cache locality,
continuous batching, memory-bandwidth floors, and two-phase
request lifecycle---then evaluated eleven schedulers spanning
load-balancing heuristics, classic and metaheuristic DAG
schedulers, a PerLLM-style constrained-bandit
(CS-UCB~\cite{yang2024perllm}), a multi-objective NSGA-II
search~\cite{yu2025llmrouting}, and three RL variants (PPO
with MLP, A3C-R2N2~\cite{tuli2022a3cr2n} with GRU encoder,
and PPO with a Decima-inspired~\cite{mao2019decima} GNN
encoder).

Our AIGC-aware PPO scheduler achieves \textbf{Pareto
dominance on mean performance} over ten strong baselines on
the (\slo{}
attainment, energy per token) plane: $+28\%$ \slo{} and
$-7\%$ energy at the canonical edge${=}5$, $\lambda{=}2$
configuration relative to the strongest single-objective
baseline, significant at $p<0.05$ on \slo{} against all
five single-objective comparators. The strongest baseline, a
knee-point NSGA-II search, is matched at canonical load,
beaten significantly in the moderate-load sweet spot
($p<0.01$), pays more energy per token in every
configuration, and spends $5.4\times$ more scheduling-time
compute. The Pareto
dominance extends robustly across edge counts $\{3, 5, 7\}$
and survives the strict A3C-R2N2 controlled comparison that
isolates the AIGC contribution from general RL machinery.

A 12-component ablation, however, reveals that \textbf{AIGC-aware
scheduling is a holistic system contribution rather than a
clever reward trick}: only continuous batching simulation
($-83\%$ \slo{} ablated) and adequate pretraining ($-38\%$
\slo{} ablated) emerge as significant components. All
individual AIGC reward and state components are statistically
indistinguishable from the Full RL baseline in isolation, and
retraining under four alternative reward-weight vectors---from
latency-first to AIGC-first---leaves both objectives
statistically unchanged while preserving Pareto dominance over
the strongest baseline.
This finding cautions the community against claiming gains
from specific AIGC reward terms without joint ablation.

We hope this work supports two longer-term directions: more
\textbf{realistic AIGC scheduling benchmarks} built on the
open platform we release, and a \textbf{research culture of
aggressive joint ablation} for AIGC scheduling claims. The
fast-moving LLM-systems landscape will reward both.
